% Edge clique interdiction: the model, Mantel, partitions, Brouwer, the
% part-size reduction, and the three open instances.
\section{Clique interdiction proofs}\label{app:clique}

\subsection{The model}

BOBILib's instance \texttt{normalClique$n$-$d$-$i$} has a graph $G=(V,E)$ with
$n=|V|$, $m=|E|$, and budget $k=\lfloor m/4\rfloor$. All variables are binary.
The leader has $w_e$ for $e\in E$ and the single row $\sum_ew_e\le k$. The
follower has $x_v$ ($v\in V$) and $y_e$ ($e\in E$) and the rows
\[
  x_u+x_v\le 1\ \ (uv\notin E),\qquad
  x_u+x_v-y_{uv}\le 1\ \ (uv\in E),\qquad
  w_{uv}+y_{uv}\le 1\ \ (uv\in E).
\]
The leader minimises $\sum_vx_v$ and the follower minimises $-\sum_vx_v$. (Every row,
coefficient and objective of the 60 files was checked against this pattern,
once by the instance reader of the clique scripts and once by the parser of a
separate review session, which classifies rows by their coefficients.)

\begin{lemma}[Follower's problem]\label{lem:clique-model}
Fix the leader's choice and let $D=\{e: w_e=1\}$. A vector $x\in\{0,1\}^V$
extends to a feasible follower point if and only if $\{v: x_v=1\}$ is a clique
of $G-D$. Hence the follower's optimum is $\omega(G-D)$, the game is zero-sum,
and the bilevel optimum is $\min\{\omega(G-D): D\subseteq E,\ |D|\le k\}$.
\end{lemma}
\begin{proof}
Let $x$ be part of a feasible point and $x_u=x_v=1$. The first family of rows
forces $uv\in E$. If $uv\in D$ then $w_{uv}=1$ forces $y_{uv}=0$, and the
second family gives $x_u+x_v\le 1$, a contradiction. So $uv\in E\setminus D$.
Conversely, for a clique of $G-D$ set $y_{uv}=1-w_{uv}$: the third family holds
with equality; the second holds for $uv\notin D$ because $y_{uv}=1$, and for
$uv\in D$ because a clique of $G-D$ does not contain both $u$ and $v$; the
first holds because a clique contains no non-edge of $G$. The leader's
objective is minus the follower's.
\end{proof}
\noindent Not formalised (the Lean file states the bounds directly in terms of
$\omega(G-D)$).

\subsection{Lower bound 3: Mantel's theorem}

\begin{theorem}[Mantel 1907, deletion form]\label{thm:mantel}
A graph on $n$ vertices with more than $\lfloor n^2/4\rfloor$ edges contains a
triangle. Consequently, if $m-k>\lfloor n^2/4\rfloor$, then $G-D$ contains a
triangle for every set $D$ of at most $k$ edges, and the bilevel optimum is
at least~$3$.
\end{theorem}
\begin{proof}
Mantel's theorem\cite{mantel1907problem} is Tur\'an's theorem\cite{turan1941extremal}
for $K_3$. After deleting at most $k$
edges at least $m-k$ edges remain, and $m-k>\lfloor n^2/4\rfloor$.
\end{proof}
{\raggedright\noindent Formalised: \texttt{Bobi/Clique.lean}, theorems \texttt{mantel},
\texttt{mantel'} and \texttt{triangle\_survives}; the instance arithmetic in
\texttt{mantel\_40}, \texttt{mantel\_50}, \texttt{mantel\_60},
\texttt{clique40}, \texttt{clique50} and \texttt{clique60}.\par}

All 30 density-0.7 instances have $(n,m,k)\in\{(40,546,136),(50,857,214),(60,1239,309)\}$,
so $m-k=410,643,930$ exceeds $\lfloor n^2/4\rfloor=400,625,900$. The density-0.9
instances satisfy the same inequality.

\subsection{Upper bounds: partitions}

\begin{proposition}[Partition upper bound]\label{prop:clique-partition}
Let $f:V\to\{1,\dots,r\}$ and let $D$ contain every edge of $G$ whose ends have
the same colour. Then $G-D$ has no clique on $r+1$ vertices. If at most $k$
edges of $G$ lie inside the colour classes, deleting exactly those edges is a
leader choice, and the bilevel optimum is at most~$r$.
\end{proposition}
\begin{proof}
$f$ is a proper colouring of $G-D$, and a clique on $r+1$ vertices needs
$r+1$ colours.
\end{proof}
{\raggedright\noindent Formalised: \texttt{Bobi/Clique.lean}, theorem
\texttt{cliqueFree\_of\_partition}.\par}

For a partition $\mathcal{P}$ of $V$ into $r$ parts (parts may be empty) let
$\mathrm{in}_G(\mathcal{P})$ be the number of edges of $G$ inside parts, and let
$w_r(G)=\min_{\mathcal{P}}\mathrm{in}_G(\mathcal{P})$. Each solution file deletes the inside edges
of a 3- or 4-partition found by tabu search and states a maximum clique of the
remaining graph, computed exactly (the review re-checks it with two exact
clique searches\cite{bron1973algorithm,tomita2006worst,carraghan1990exact});
the files pass the checks of Section~\ref{app:verification}.

\subsection{Lower bound 4: Brouwer's theorem}

Let $t_r(n)$ be the number of edges of the balanced complete $r$-partite graph
on $n$ vertices (Tur\'an's number), and
$T(n,r)=\binom n2-t_r(n)$, the number of edges of the disjoint union of $r$
cliques of balanced sizes.

\begin{theorem}[Brouwer 1981, Theorem~2\cite{brouwer1981lotto}; quoted, with his $G$ renamed $H$ and his $m:=\lfloor n/t\rfloor$ written out]\label{thm:brouwer}
Let $H$ be a graph on $n$ points such that
$\alpha(H)\le t$, and $n>2t$. If $H$ has no more than $T(n,t)+\lfloor n/t\rfloor-2$
edges, then $V(H)$ can be written as a union of $t$ cliques. Moreover, this
bound is best possible.
\end{theorem}

Here $\alpha(H)$ is the independence number. The theorem is restated by
Tyomkyn and Uzzell\cite{tyomkyn2015strong} (their equation~(1) and
Theorem~1.1), who credit Brouwer with the first explicit proof, and, in the
complemented form below for $t\ge 2$, by Kor\'andi, Roberts and
Scott\cite{korandi2021exact} (their Theorem~1.1). Brouwer's report is not
refereed, but the statement is: Tyomkyn and Uzzell give a new proof of it in
their Section~2, Hanson and Toft and Kang and Pikhurko found it
again,\cite{hanson1991ksaturated,kang2005maximum} and Amin, Faudree, Gould and
Sidorowicz\cite{amin2013non} determine the largest number of edges of a maximal
$K_p$-free graph that is not $(p-1)$-partite, which gives the same threshold
(a $K_p$-free graph that is not $(p-1)$-partite stays so when edges are added
until it is maximal). We use the theorem in complemented form.

\begin{corollary}[Complement form]\label{cor:brouwer}
Let $t\ge 1$ and $n\ge 2t+1$. Every $K_{t+1}$-free graph $F$ on $n$ vertices with
$e(F)\ge t_t(n)-\lfloor n/t\rfloor+2$ is $t$-partite (its vertex set is the
union of $t$ independent sets, some possibly empty).
\end{corollary}
\begin{proof}
Apply Theorem~\ref{thm:brouwer} to the complement $H=\overline F$:
$\alpha(H)=\omega(F)\le t$, $n>2t$, and
$e(H)=\binom n2-e(F)\le\binom n2-t_t(n)+\lfloor n/t\rfloor-2=T(n,t)+\lfloor n/t\rfloor-2$.
So $V$ is a union of $t$ cliques of $H$, that is, of $t$ independent sets of
$F$. Assigning each vertex to one of the sets containing it gives a partition
into independent sets.
\end{proof}
\noindent Not formalised: Brouwer's theorem is not in Mathlib, and it is
neither proved nor assumed in the Lean files. The statement above was read
in the scanned report, and the translation was checked by computer (for
$t=3$: $T(n,3)+t_3(n)=\binom n2$).

\begin{theorem}[Density 0.9: value at least 4]\label{thm:clique-lb4}
Let $n\ge 7$, $m-k\ge t_3(n)-\lfloor n/3\rfloor+2$ and $w_3(G)>k$. Then $G-D$
contains a $K_4$ for every $D\subseteq E$ with $|D|\le k$, so the bilevel optimum is
at least~$4$.
\end{theorem}
\begin{proof}
Suppose $|D|\le k$ and $G-D$ is $K_4$-free. It has $m-|D|\ge m-k$ edges, so by
Corollary~\ref{cor:brouwer} with $t=3$ ($n\ge 7=2t+1$) it is 3-partite with some
partition $\mathcal{P}$. No edge of $G-D$ lies inside a part, so every edge of $G$
inside a part lies in $D$, and $|D|\ge\mathrm{in}_G(\mathcal{P})\ge w_3(G)>k$, a
contradiction.
\end{proof}
\noindent Not formalised (it rests on Theorem~\ref{thm:brouwer}).

For the density-0.9 instances $(n,m,k)\in\{(40,702,175),(50,1102,275),(60,1593,398)\}$,
so $m-k=527,827,1195$ and $t_3(n)-\lfloor n/3\rfloor+2=522,819,1182$
($t_3(n)=533,833,1200$): the hypothesis holds with margins $5$, $8$ and $13$,
while $m-k<t_3(n)$, so Tur\'an's theorem alone does not give a $K_4$. It
remains to show $w_3(G)>k$.

\subsection{Reducing $w_3(G)>k$ to finitely many part-size cases}

Let $\overline E$ be the edge set of the complement $\overline G$, so
$|\overline E|=\binom n2-m$, and for a partition $\mathcal{P}$ let
$\mathrm{cut}_{\overline G}(\mathcal{P})$ be the number of edges of $\overline G$ joining
different parts.

\begin{proposition}[Part-size reduction]\label{prop:clique-sizes}
For a partition $\mathcal{P}=(P_1,P_2,P_3)$ of $V$ with part sizes $s_1,s_2,s_3$,
\[
  \mathrm{in}_G(\mathcal{P})=\sum_{i=1}^3\binom{s_i}2-|\overline E|+\mathrm{cut}_{\overline G}(\mathcal{P}).
\]
Hence, with $b(s)=k+|\overline E|-\sum_i\binom{s_i}2$, $w_3(G)>k$ holds if and only
if, for every size vector $s_1\ge s_2\ge s_3\ge 0$ with $s_1+s_2+s_3=n$ and
$b(s)\ge 0$, no partition with these part sizes has
$\mathrm{cut}_{\overline G}(\mathcal{P})\le b(s)$.
\end{proposition}
\begin{proof}
Each of the $\sum_i\binom{s_i}2$ pairs inside a part is an edge of $G$ or of
$\overline G$, and the edges of $\overline G$ inside parts number
$|\overline E|-\mathrm{cut}_{\overline G}(\mathcal{P})$. So $\mathrm{in}_G(\mathcal{P})\le k$ iff
$\mathrm{cut}_{\overline G}(\mathcal{P})\le b(s)$. Relabelling the parts, every partition
has sorted sizes, and since $\mathrm{cut}_{\overline G}\ge 0$, size vectors with
$b(s)<0$ admit no partition with $\mathrm{in}_G\le k$.
\end{proof}
\noindent Not formalised (checked in exact integer arithmetic by script, and
on random partitions of every instance).

Since $n$, $m$ and $k$ are the same for all ten graphs of a group, so are the
size vectors with $b(s)\ge 0$:
\begin{center}
\small
\begin{tabular}{lll}
\hline
$n$ & $|\overline E|$ & size vectors $s$ with $b(s)$ in brackets\\
\hline
40 & 78 & (14,13,13)\,[6], (14,14,12)\,[5], (15,13,12)\,[4], (15,14,11)\,[2], (16,12,12)\,[1], (16,13,11)\,[0]\\
50 & 123 & (17,17,16)\,[6], (18,16,16)\,[5], (18,17,15)\,[4], (19,16,15)\,[2], (18,18,14)\,[1], (19,17,14)\,[0]\\
60 & 177 & (20,20,20)\,[5], (21,20,19)\,[4], (21,21,18)\,[2], (22,19,19)\,[2], (22,20,18)\,[1]\\
\hline
\end{tabular}
\end{center}
That is $10\cdot 6+10\cdot 6+10\cdot 5=170$ cases (instance, size vector), each
asking for a partition of $V$ with prescribed part sizes and at most $b(s)\le 6$
edges of the sparse complement between parts. All 170 were decided
infeasible by five methods: CP-SAT\cite{perron2026cpsat} and
SCIP\cite{hojny2025scip} on two different models, an exact
solver-free depth-first search (tested against brute force), and, in a
separate review session with its own code, HiGHS on a third model with an explicit budget row and
a separate exact branch and bound written in C (tested against brute force
on 450 random graphs and on planted instances). The three solver verdicts are
solver-certified (class~(ii), Section~\ref{app:classes}). For an exact
certificate, each case also has a branch-and-bound tree
(\texttt{results/clique/certificates/}), checked by the integer checker
\texttt{scripts/check\_clique\_certificates.py}: it reads the instance files
itself, re-derives the graph, $n$, $m$ and $k$ from the rows, checks the
hypothesis $m-k\ge t_3(n)-\lfloor n/3\rfloor+2$, enumerates the size vectors
with $b(s)\ge 0$, and verifies at every leaf of each tree a lower bound on
$\mathrm{cut}_{\overline G}$ above $b(s)$, using only Python's integers and
no project code (\CliqueCheckerLines{} lines of code). It passes on all
\CliqueCertCases{} cases of the \CliqueCertInstances{} instances
(\texttt{check.log}): the trees have \CliqueCertNodes{} nodes and
\CliqueCertLeaves{} leaves in all, the largest \CliqueCertMaxNodes{} nodes, and
the check takes about \CliqueCertSeconds{}~s. With Theorem~\ref{thm:clique-lb4} and
Proposition~\ref{prop:clique-partition} (4-partitions with at most $k$ inside
edges), all 30 density-0.9 instances have value exactly~$4$: an exact
certificate (class~(i)) resting on Brouwer's published theorem, which is
neither formalised nor reproved here. The case computations are not
formalised proofs.

\subsection{The three open instances}

For normalClique60-70-2, -3 and -5 ($n=60$, $m=1239$, $k=309$) Mantel gives
value $\ge 3$ and a 4-partition gives value $\le 4$. Let $\tau_4(G)$ be the
fewest edges meeting every $K_4$ of $G$. By Lemma~\ref{lem:clique-model} the
value is $3$ if and only if $\tau_4(G)\le k$. Deleting the inside edges of a
3-partition destroys every $K_4$ (two of any four vertices share a part), so
$\tau_4(G)\le w_3(G)$, and $w_3(G)\le k$ would give value~$3$; the best
3-partitions found have $311$, $310$ and $310$ inside edges, against $k=309$.

The converse argument of Theorem~\ref{thm:clique-lb4} is not available
here. It would need every $K_4$-free subgraph with at least $m-k=930$ edges
to be 3-partite, but $930<t_3(60)-20+2=1182$, so Corollary~\ref{cor:brouwer}
does not apply, and its conclusion is false at this edge count: the disjoint
union of the wheel $W_5$ (a 5-cycle plus a vertex joined to all of it; 6
vertices, 10 edges, $K_4$-free, chromatic number~4) and any 3-partite graph
with $920$ edges on the other 54 vertices ($920\le t_3(54)=972$) is a
$K_4$-free graph on 60 vertices with 930 edges that is not 3-partite. So a
proof of value~4 must bound $\tau_4(G)$ itself from below. Linear
programming bounds on $\tau_4$ (rows for every $K_4$, for maximal cliques and
for vertex sets whose induced subgraph must lose a Tur\'an excess of edges)
stop near $296$, below the $310$ needed, and searches for a transversal with
at most $309$ edges that are not restricted to 3-partite remainders stop at
the best 3-partition found. The three instances stay open with
$3\le\text{value}\le 4$.
